% Mapa mental do Subcomponente 3.2.b
\begin{tikzpicture}[
    ramo/.style={draw=#1, fill=#1!10, rounded corners=3pt, align=left,
                 text width=4.5cm, inner sep=5pt, font=\scriptsize},
    liga/.style={draw=#1, line width=1.3pt}
  ]
  \colorlet{c1}{azulpsh}
  \colorlet{c2}{green!45!black}
  \colorlet{c3}{orange!85!black}
  \colorlet{c4}{violet!75!black}
  \colorlet{c5}{teal!80!black}
  \colorlet{c6}{cyan!60!black}
  \colorlet{c7}{brown!80!black}
  \colorlet{c8}{red!60!black}

  % --- lado esquerdo ------------------------------------------------
  \node[ramo=c1, anchor=north] (m1) at (-5.9,0) {\textbf{\color{c1}Metas}\\[2pt]
    \textbullet\ 83 sistemas coletivos de água (IAT)\\
    \textbullet\ 222 sistemas pelo Sanepar Rural\\
    \textbullet\ esgoto em 4.067 domicílios das comunidades\\
    \textbullet\ esgoto em 2.300 domicílios em mananciais};
  \node[ramo=c2, below=0.45cm of m1] (m2) {\textbf{\color{c2}Seleção e mapa (1.2 e 1.3)}\\[2pt]
    \textbullet\ $\geq$ 70 domicílios por comunidade\\
    \textbullet\ IPDM do município\\
    \textbullet\ associações de moradores existentes\\
    \textbullet\ mananciais prioritários do PSH-PR};
  \node[ramo=c3, below=0.45cm of m2] (m3) {\textbf{\color{c3}Responsáveis (1.4)}\\[2pt]
    \textbullet\ IAT: executor\\
    \textbullet\ IDR-Paraná e Sanepar: coexecutores\\
    \textbullet\ SECID, prefeituras e comunidades\\
    \textbullet\ empresas contratadas e UGP};
  \node[ramo=c4, below=0.45cm of m3] (m4) {\textbf{\color{c4}Recursos e prazo (1.12 e 1.13)}\\[2pt]
    \textbullet\ R\$ 183,0 milhões do BIRD (IAT)\\
    \textbullet\ R\$ 67,1 milhões da Sanepar\\
    \textbullet\ sete anos: 2027 a 2033\\
    \textbullet\ preparação em 2026};

  % --- lado direito -------------------------------------------------
  \node[ramo=c5, anchor=north] (m5) at (5.9,0) {\textbf{\color{c5}Modelo de gestão (1.6)}\\[2pt]
    \textbullet\ gestão compartilhada: Estado, município e comunidade\\
    \textbullet\ associações regularizadas\\
    \textbullet\ trabalho técnico-social com as famílias};
  \node[ramo=c6, below=0.45cm of m5] (m6) {\textbf{\color{c6}Abastecimento de água (1.7)}\\[2pt]
    \textbullet\ poço, outorga, projeto e licenças\\
    \textbullet\ obras com supervisão e fiscalização\\
    \textbullet\ entrega à entidade gestora\\
    \textbullet\ Programa Sanepar Rural};
  \node[ramo=c7, below=0.45cm of m6] (m7) {\textbf{\color{c7}Esgotamento sanitário (1.8)}\\[2pt]
    \textbullet\ solução individual por domicílio\\
    \textbullet\ frentes comunidades e mananciais\\
    \textbullet\ verificação técnica e medição\\
    \textbullet\ pagamento após comprovação};
  \node[ramo=c8, below=0.45cm of m7] (m8) {\textbf{\color{c8}Qualidade e salvaguardas (1.10 e 1.11)}\\[2pt]
    \textbullet\ medidas ambientais e sociais em todas as fases\\
    \textbullet\ canal de queixas nas comunidades\\
    \textbullet\ evidência e validação de cada produto\\
    \textbullet\ metas e indicadores anuais};

  % --- centro (entre a 2ª e a 3ª caixa da coluna esquerda) -----------
  \node[ellipse, draw=azulpsh, fill=azulpsh, text=white, align=center,
        minimum width=4.3cm, minimum height=2.3cm, font=\small\bfseries] (c) at ($(m2.south)!0.5!(m3.north) + (5.9,0)$)
    {Subcomponente 3.2.b\\[2pt]\footnotesize\mdseries Água e esgoto\\ no meio rural};

  % --- ligações -----------------------------------------------------
  \begin{scope}[on background layer]
  \foreach \n/\cor in {m1/c1, m2/c2, m3/c3, m4/c4}
    \draw[liga=\cor] (c.west) to[out=180, in=0] (\n.east);
  \foreach \n/\cor in {m5/c5, m6/c6, m7/c7, m8/c8}
    \draw[liga=\cor] (c.east) to[out=0, in=180] (\n.west);
  \end{scope}
\end{tikzpicture}
